\documentclass[10pt,a4paper]{amsart}
\usepackage[margin=19mm]{geometry}
\usepackage{amsmath,amssymb,mathtools}
\usepackage[scr=rsfs,cal=euler]{mathalfa}
\usepackage{xfrac}
\usepackage[unicode,hidelinks,bookmarks=false]{hyperref}
\newcommand{\ie}{i.e.\ }
\newcommand{\cf}{cf.\ }
\newcommand{\resp}{respectively\ }
\newcommand{\Subsection}{\subsection}
\providecommand{\nobreakdash}{\nobreak\hbox{-}}
\setlength{\emergencystretch}{2em}
\multlinegap0pt
\usepackage{amssymb}
\newtheorem{proposition}{Proposition}
\newcommand{\LOk}{L_{\mathrm{Ok}}}
\title{Integral Shell Polytopes of Composition Algebras}
\author{Daniele Corradetti}
\date{Source: arXiv:2605.09458v1 (CC BY 4.0)}
\begin{document}
\maketitle

\noindent\emph{Source and licence.} Integral Shell Polytopes of Composition Algebras. Version arXiv:2605.09458v1 (CC BY 4.0), distributed by arXiv under the Creative Commons Attribution 4.0 International licence, http://creativecommons.org/licenses/by/4.0/, as recorded in the arXiv licence metadata for this exact version and shown on the abstract page https://arxiv.org/abs/2605.09458v1. Full licence notice: \texttt{licenses/arxiv-2605.09458-CC-BY-4.0.txt}.\par\smallskip

\noindent\textit{Editorial scope.} Counted unit: the article's proposition prop:D8-root (source0.tex lines 729--733). The article's complete original proof of it is delivered with it (source0.tex lines 735--754, one \texttt{proof} environment of 20 lines). No mathematics is written, repaired or re-derived: the statement and the proof are the article's own lines, in the article's own notation, and no retained line is substituted or re-flowed.  No further source-local context is needed: every cross-reference of every retained line resolves inside the two delivered spans.  Nothing was omitted from any retained span, so the package carries no omission record.  No retained line contains a citation, so the wrapper carries no bibliography and no bibliography entry.  No retained line contains a figure environment, an image-inclusion command or drawing code, so no image is delivered and the package declares no asset dependency.  Editorial formatting only, in addition: an AMS \texttt{amsart} preamble that restates the article's own declarations and macros used by the retained lines, namely the environments \texttt{proposition} on separate counters, together with the standard packages those lines need.  The retained environments print in this document as Proposition 1 in order of appearance, whereas the article prints the same items with the numbering of its own full document.  The complete native source of the article as distributed in the arXiv e-print for this exact version is bundled unchanged as \texttt{sources/200222/source0.tex}.
\begin{proposition}
\label{prop:D8-root}
The shell \(S_8(\LOk)\) has \(112\) vectors and is the root polytope of type
\(D_8\).
\end{proposition}

\begin{proof}
The shell is antipodal and spans rank \(8\). Numerical computation (see Section 7)
selects eight simple roots in the shell. After normalizing the shell vectors
to squared length \(2\), the Cartan matrix is
\begin{equation}
\begin{pmatrix}
2&-1&0&0&0&0&0&0\\
-1&2&-1&0&0&0&0&0\\
0&-1&2&-1&0&0&0&0\\
0&0&-1&2&-1&0&0&0\\
0&0&0&-1&2&-1&0&0\\
0&0&0&0&-1&2&-1&-1\\
0&0&0&0&0&-1&2&0\\
0&0&0&0&0&-1&0&2
\end{pmatrix}.
\end{equation}
This is a Cartan matrix of type \(D_8\). One then verifies
that the reflections generated by these simple roots preserve the shell and
produce an orbit of \(112\) vectors, equal to all of \(S_8(\LOk)\).
\end{proof}

\end{document}
